\documentclass[11pt,letterpaper]{article}
\usepackage[T1]{fontenc}
\usepackage[utf8]{inputenc}
\usepackage{amsmath,amsthm}
\usepackage{newpxtext,newpxmath}
\usepackage[margin=1in,headheight=16pt,footskip=30pt]{geometry}
\usepackage{microtype}
\usepackage{mathtools}
\usepackage{xcolor}
\usepackage{booktabs,tabularx,array}
\usepackage{enumitem}
\usepackage{titlesec,fancyhdr}
\usepackage[most]{tcolorbox}
\usepackage{listings}
\usepackage[numbers,sort&compress]{natbib}
\usepackage{xurl}
\usepackage[colorlinks=true,linkcolor=OliveDark,citecolor=OliveDark,urlcolor=OliveDark,
  pdftitle={A Square-Root Obstruction to Evaluating Increasing Hahn Series at Omega},
  pdfauthor={Research note prepared for Vladimir Reshetnikov},
  pdfsubject={Negative answer to Lipparini Problem 7.7; Hahn series and surreal numbers;
    nonnegative-coefficient semiring strengthening},
  pdfkeywords={surreal numbers, Hahn series, ordered fields, formal power series,
    semirings, well-quasi-orders, ordinal-indexed supports}]{hyperref}
\usepackage[nameinlink,noabbrev]{cleveref}
\definecolor{OliveDark}{HTML}{505A3C}
\definecolor{OlivePale}{HTML}{F0F2E9}
\definecolor{Ink}{HTML}{222820}
\definecolor{Quiet}{HTML}{66685E}
\definecolor{Rule}{HTML}{B6BCA7}
\definecolor{Warm}{HTML}{FAF8F2}
\color{Ink}
\setlength{\parindent}{1.2em}
\setlength{\parskip}{0.28em}
\setlength{\emergencystretch}{2em}
\linespread{1.055}
\titleformat{\section}{\Large\bfseries\color{OliveDark}}{\thesection}{0.7em}{}
\titleformat{\subsection}{\large\bfseries\color{OliveDark}}{\thesubsection}{0.7em}{}
\titleformat{\subsubsection}{\normalsize\bfseries}{\thesubsubsection}{0.6em}{}
\titlespacing*{\section}{0pt}{2.1ex plus 1ex minus .2ex}{1.0ex plus .2ex}
\titlespacing*{\subsection}{0pt}{1.8ex plus 1ex minus .2ex}{.7ex plus .2ex}
\pagestyle{fancy}
\fancyhf{}
\fancyhead[L]{\small\color{Quiet}Increasing Hahn series and surreal evaluation}
\fancyhead[R]{\small\color{Quiet}Research note}
\fancyfoot[C]{\small\thepage}
\renewcommand{\headrulewidth}{0.35pt}
\renewcommand{\footrulewidth}{0pt}
\newtheorem{theorem}{Theorem}[section]
\newtheorem{propositionx}[theorem]{Proposition}
\newtheorem{lemmax}[theorem]{Lemma}
\newtheorem{corollaryx}[theorem]{Corollary}
\theoremstyle{definition}
\newtheorem{definitionx}[theorem]{Definition}
\newtheorem{examplex}[theorem]{Example}
\theoremstyle{remark}
\newtheorem{remarkx}[theorem]{Remark}
% All of these share the theorem counter, so cleveref would otherwise call
% every one of them a ``theorem''. Aliasing the counter inside each
% environment gives the right name while keeping one running number.
\crefname{lemma}{Lemma}{Lemmas}
\Crefname{lemma}{Lemma}{Lemmas}
\crefname{proposition}{Proposition}{Propositions}
\Crefname{proposition}{Proposition}{Propositions}
\crefname{corollary}{Corollary}{Corollaries}
\Crefname{corollary}{Corollary}{Corollaries}
\crefname{definition}{Definition}{Definitions}
\Crefname{definition}{Definition}{Definitions}
\crefname{example}{Example}{Examples}
\Crefname{example}{Example}{Examples}
\crefname{remark}{Remark}{Remarks}
\Crefname{remark}{Remark}{Remarks}
\newenvironment{lemma}{\crefalias{theorem}{lemma}\lemmax}{\endlemmax}
\newenvironment{proposition}{\crefalias{theorem}{proposition}\propositionx}{\endpropositionx}
\newenvironment{corollary}{\crefalias{theorem}{corollary}\corollaryx}{\endcorollaryx}
\newenvironment{definition}{\crefalias{theorem}{definition}\definitionx}{\enddefinitionx}
\newenvironment{example}{\crefalias{theorem}{example}\examplex}{\endexamplex}
\newenvironment{remark}{\crefalias{theorem}{remark}\remarkx}{\endremarkx}
\newcommand{\No}{\mathbf{No}}
\newcommand{\R}{\mathbb R}
\newcommand{\Q}{\mathbb Q}
\newcommand{\N}{\mathbb N}
\newcommand{\Z}{\mathbb Z}
\newcommand{\Hu}[1]{\mathcal H_{\uparrow}(#1)}
\newcommand{\Hfull}{\mathcal H_{\uparrow}(\No)}
\newcommand{\supp}{\operatorname{supp}}
\newcommand{\lc}{\operatorname{lc}}
\newcommand{\ev}{\operatorname{ev}}
\newcommand{\ord}{\operatorname{ord}}
\newcommand{\vexp}{\operatorname{v}}
\newcommand{\Inf}{\mathfrak I}
\newcommand{\degree}{\operatorname{deg}}
\newcommand{\coeff}[2]{[#1^{#2}]}
\newcommand{\Qnn}{\Q_{\ge0}}
\newcommand{\seqsum}{\mathop{\sum}\nolimits^{\mathrm H}}
\newtcolorbox{resultbox}{enhanced,breakable,colback=OlivePale,colframe=Rule,
  boxrule=.5pt,arc=1.5mm,left=3mm,right=3mm,top=2mm,bottom=2mm}
\lstset{basicstyle=\small\ttfamily,columns=fullflexible,keepspaces=true,
  breaklines=true,frame=single,rulecolor=\color{Rule},backgroundcolor=\color{Warm},
  xleftmargin=.3em,xrightmargin=.3em,showstringspaces=false}
\setlist{itemsep=.15em,topsep=.4em}
\setcounter{tocdepth}{1}
\hypersetup{bookmarksdepth=2}
\begin{document}
\hypersetup{pageanchor=false}%
\begin{titlepage}
\thispagestyle{empty}
\vspace*{0.05in}
{\color{OliveDark}\small\sffamily ORDER THEORY \; / \; HAHN FIELDS \; / \; SURREAL NUMBERS\par}
\vspace{0.28in}
{\Huge\bfseries A square-root obstruction\par}
\vspace{.12in}
{\LARGE\bfseries to evaluating increasing\par Hahn series at $\omega$\par}
\vspace{.28in}
{\large A negative answer to Lipparini's Problem 7.7,\par
with a nonnegative-coefficient semiring strengthening,\par
infinitesimal evaluation and monomial classification\par}
\vspace{.22in}
{\color{Quiet}20 September 2026\par}
\vspace{.2in}
\begin{resultbox}
The decisive identity already lies in $\Q[[x]]$:
\[
 B(x)=1-\frac{x}{2}-\frac{x^2}{8}-\frac{x^3}{16}-\frac{5x^4}{128}-\cdots,
 \qquad B(x)^2=1-x.
\]
A ring homomorphism sending $x$ to $\omega$ would give
\[
 \Phi(B)^2=1-\omega<0,
\]
which is impossible in the ordered field of surreal numbers.
Subtraction is not to blame: the all-positive pair
$H(x)=\sum_{n\ge0}\binom{2n}{n}4^{-n}x^n$ and $G(x)=\sum_{n\ge0}x^n$
satisfies $H^2=G$ and $G=1+xG$, so even a homomorphism of the
\emph{semiring} $\Qnn[[x]]$ cannot send $x$ to any $u\ge1$.
\end{resultbox}
\vspace{.25in}
\noindent\textbf{Abstract.}
We give a negative answer to the increasing-support Hahn-series evaluation problem
posed as Problem 7.7 in the April 2026 extended version of Paolo Lipparini's paper
on monotone infinitary ordinal operations. The obstruction is elementary and
algebraic: the formal square root of $1-x$ exists in the source, whereas $1-\omega$
has no square root in an ordered field. No continuity, order preservation, or
preservation of infinite sums is assumed. We then strengthen the obstruction
along two axes at once: it survives on the semiring $\Qnn[[x]]$ of
nonnegative-coefficient series, where there are no additive inverses at all,
and there it already excludes every value $u\ge1$. We prove that the possible
images of $x$ under unital homomorphisms $\Q[[x]]\to\No$ are exactly the surreal
infinitesimals, including zero. For general Hahn series we classify the pure
monomial prescriptions that admit a homomorphism: the exponent map must be
additive and strictly order-reversing, and this condition is sufficient. We also
explain the classical exponent-reversing isomorphism, isolate the obstruction
inside a countable finitely generated algebra, and supply self-contained
summability arguments and exact finite regression checks.

\vspace{.16in}
\noindent\textbf{Provenance.}
This report is a merge of two independently produced research reports on the
same question, circulated as \texttt{hahn-evaluation-at-omega} and
\texttt{reversed-hahn-series}. Both reached the same negative answer by the
same argument: each constructed $B$ from the identical quadratic recursion,
each verified $B^2=1-x$ coefficient by coefficient, and each applied the
hypothetical homomorphism only to that one finite identity. That shared core
theorem is therefore proved exactly once below, in \cref{lem:witness} and
\cref{thm:main}. The supporting material genuinely diverged, and this document
is the union of the two toolkits rather than a deduplication of them: the
semiring strengthening of \cref{sec:semiring}, the general unit square root
of \cref{lem:unit-root}, the elementary two-support argument of
\cref{lem:two-supports}, and several scope remarks in \cref{sec:audit} come
from the second report, while the three-way evaluation criterion, the
order-forcing theorem, the monomial classification over an arbitrary exponent
group, and the appendix on Higman's lemma come from the first. Where the two
reports phrased an attribution or a priority caveat differently, the more
cautious wording has been kept.
\vfill
\noindent{\small\color{Quiet}
Independent research note prepared for Vladimir Reshetnikov. The application to
the stated problem is proved here; the Hahn-field background is classical.
This manuscript has not undergone independent refereeing or proof-assistant verification.}
\end{titlepage}
\hypersetup{pageanchor=true}%

\tableofcontents
\vspace{1.0em}
\begin{resultbox}
\textbf{Reading guide.} Section~\ref{sec:main} contains the complete negative
solution and does not depend on the later sections.
Section~\ref{sec:semiring} is an independent second obstruction using only
nonnegative coefficients and semiring operations; it needs nothing from
Section~\ref{sec:main} and is not implied by it. Sections~\ref{sec:inf}
and~\ref{sec:classification} describe exactly what can be evaluated instead.
Section~\ref{sec:audit} collects the scope limits and the misreadings the
result invites. Appendix~\ref{app:summability} supplies the order-combinatorial
details used only in the extensions. The supplied code checks finite identities;
it is not a substitute for the proofs.
\end{resultbox}
\clearpage

\section{The problem and its precise scope}\label{sec:problem}

\subsection{An explicitly posed evaluation question}
Lipparini's extended paper \citep[Problem 7.7, p.~38]{Lipparini2026} considers
formal sums
\begin{equation}\label{eq:source}
 s=\sum_{\beta<\alpha} r_\beta x^{a_\beta},
 \qquad r_\beta\in\R\setminus\{0\},
 \qquad \gamma<\beta<\alpha\Longrightarrow a_\gamma<a_\beta,
\end{equation}
where $\alpha$ is an ordinal and the exponents are surreal numbers. The question
is whether there is a ring homomorphism into $\No$ satisfying
\begin{equation}\label{eq:prescription}
 \Phi(x^a)=\omega^a\qquad(a\in\No).
\end{equation}
The support direction is essential: the source exponents increase, whereas
Conway normal forms have decreasing exponents. The stated question asks for
ordinary ring operations to be preserved, not necessarily infinite sums.

The formula $a=0$ in \eqref{eq:prescription} forces $\Phi(1)=1$, regardless of
whether one's convention includes unitality in the term ``ring homomorphism.''
The formula $a=1$ forces $\Phi(x)=\omega$. These are the only two monomial
requirements needed for the negative answer.

\subsection{Notation and arithmetic conventions}
Two warnings are worth stating before any notation. First, the addition and
multiplication used throughout this note --- in the target and in the exponent
group alike --- are the surreal \emph{ordered-field} operations, never the
noncommutative ordinal operations; $1-\omega$ below is a negative surreal
number, not an attempted ordinal subtraction. Second, we do not assume that
$\Phi$ fixes every real number. Taking $a=0$ in \eqref{eq:prescription}
already forces $\Phi(1)=1$, and a unital ring homomorphism necessarily fixes
$\Q$. Rational coefficients are all that the obstruction uses, so the status
of arbitrary real coefficients never arises in the negative proof.

For a linearly ordered abelian group $G$, written additively, let
\[
 \Hu{G}=\left\{\sum_{a\in A} r_a x^a:
  A\subseteq G\text{ is well-ordered},\ r_a\in\R\right\}.
\]
Zero coefficients are discarded. The zero series has empty support. Each
support is a set, including when $G=\No$ is a proper class. Addition is
coefficientwise, and multiplication is the Hahn convolution
\begin{equation}\label{eq:convolution}
 (st)_c=\sum_{a+b=c}s_a t_b.
\end{equation}
For fixed supports and fixed $c$, only finitely many summands are nonzero.
We give the relevant well-ordering proof in Appendix~\ref{app:summability}.
The source of the problem is $\Hfull$.

All exponent addition in \eqref{eq:convolution}, when $G=\No$, is surreal
addition. All target addition and multiplication are surreal field operations.
In particular, $1-\omega$ is a negative surreal number; it is not an expression
in ordinal subtraction. We use $\omega^a$ for Conway's omega-map, which satisfies
\begin{equation}\label{eq:omega}
 \omega^0=1,\qquad \omega^{a+b}=\omega^a\omega^b,
 \qquad a<b\Longrightarrow \omega^a<\omega^b.
\end{equation}
This notation is not a definition of $\exp(a\log\omega)$ for arbitrary surreal
$a$. No exponential-function compatibility is imposed in this note.

\subsection{Which background results are being used?}
Conway normal forms identify $\No$ with real-coefficient generalized series
whose surreal exponents are reverse well-ordered. Each such series is a unique
surreal number, and its field operations are coefficientwise addition and
convolution. The sign is the sign of its leading coefficient. A pure monomial
$\omega^b$ is infinitesimal exactly when $b<0$. These facts, together with
\eqref{eq:omega}, can be found in the primary treatments
\citep[pp.~175--176]{DriesEhrlich2001} and
\citep[Sections 2.3--2.6]{BerarducciMantova2018}; the same conventions are
recorded in \citep[Sections 1 and 2]{KKS} and \citep{BG}.

For the main contradiction, only the much weaker facts that $\No$ is an
ordered field and $\omega>1$ are required. None of the detailed normal-form
theory is needed to construct the obstructing source series.

\subsection{What is, and is not, being claimed}\label{sec:claims}
The question is still presented as an open problem in the version of the
source inspected for this note, and the argument below gives a complete
negative answer to that formulation. The square-root construction and the
Hahn-field methods themselves are classical; what is claimed here is their
application to the specified question, together with an explicit account of
the boundary cases. No exhaustive priority determination has been made and
this manuscript has not been independently refereed.

The separate problem of constructing useful infinitary operations on
arbitrary surreal sequences or on combinatorial games is not settled here.
In particular, Lipparini's broader Problem 7.5 is \emph{not} refuted by the
failure of this particular ring-homomorphism approach. A partial operation,
an operation defined on a smaller domain, or an operation that is not
required to be multiplicative can each evade the present obstruction.
\Cref{sec:audit} returns to these limits in detail.

\section{The negative answer}\label{sec:main}

\subsection{Constructing the witness without analytic convergence}
An element of $\Q[[x]]$ is a coefficient sequence $(c_n)_{n\ge0}$, written
$\sum_{n\ge0}c_nx^n$; equality of formal series means equality of every
coefficient, and no numerical substitution, convergence, or limiting process
is part of the definition. The product of $\sum c_ix^i$ and $\sum d_jx^j$ has
$\coeff{x}{n}$ equal to the \emph{finite} sum $\sum_{i=0}^{n}c_id_{n-i}$.

The map
\begin{equation}\label{eq:ordinary-embedding}
 \Q[[x]]\longrightarrow\Hfull,
 \qquad \sum_{n\ge0}c_nx^n\longmapsto\sum_{n\ge0}c_nx^n,
\end{equation}
is therefore an \emph{injective ring homomorphism}. Once zero coefficients are
discarded, the support is a subset of $\{0,1,2,\ldots\}$ and hence a
well-ordered set of surreal exponents; the Hahn convolution
\eqref{eq:convolution} restricts on such supports to precisely the displayed
finite Cauchy convolution; and distinct coefficient sequences give distinct
supports-with-coefficients. We use \eqref{eq:ordinary-embedding} silently from
now on and regard $\Q[[x]]$ as a subring of $\Hfull$.

\begin{lemma}[The formal square-root witness]\label{lem:witness}
There is a \emph{unique} series $B\in\Q[[x]]$ with constant coefficient $1$
such that
\begin{equation}\label{eq:witness}
 B^2=1-x.
\end{equation}
Its support is exactly $\N$, and all its nonconstant coefficients are negative
dyadic rationals.
\end{lemma}
\begin{proof}
Define rational coefficients recursively by
\begin{equation}\label{eq:recurrence}
 b_0=1,\qquad b_1=-\frac12,\qquad
 b_n=-\frac12\sum_{k=1}^{n-1}b_kb_{n-k}\quad(n\ge2),
\end{equation}
and set $B=\sum_{n\ge0}b_nx^n$. This is a formal series: it is defined by its
coefficient sequence, without a convergence requirement.

The constant coefficient of $B^2$ is $1$, and its coefficient of $x$ is
$2b_1=-1$. For $n\ge2$ the coefficient of $x^n$ is
\[
 \sum_{k=0}^{n}b_kb_{n-k}
 =2b_n+\sum_{k=1}^{n-1}b_kb_{n-k}=0.
\]
This proves \eqref{eq:witness}. Read in the other direction, the same
computation shows that \eqref{eq:recurrence} is forced: with $b_0=1$ fixed,
requiring each coefficient of $B^2$ to match that of $1-x$ determines exactly
one new rational number at every stage. The recursion therefore proves
existence and uniqueness together.

The recursion preserves dyadic rationals.
Moreover, $b_1<0$; if $b_1,\ldots,b_{n-1}$ are negative, every summand in the
recursion for $b_n$ is positive, so $b_n<0$. Induction proves that none of the
coefficients vanishes.
\end{proof}

\begin{theorem}[Negative answer to Problem 7.7]\label{thm:main}
There is no ring homomorphism
\[
 \Phi:\Hfull\longrightarrow\No
\]
satisfying $\Phi(x^a)=\omega^a$ for every surreal $a$. More strongly, there is
no unital homomorphism $\Q[[x]]\to F$ into an ordered field $F$ that sends $x$
to an element greater than $1$.
\end{theorem}
\begin{proof}
Let $u$ be the image of $x$ and use the series in \cref{lem:witness}. A unital
ring homomorphism would imply
\[
 \Phi(B)^2=\Phi(B^2)=\Phi(1-x)=1-u.
\]
If $u>1$, the right-hand side is negative, whereas a square in an ordered
field is nonnegative. In the problem, $u=\omega>1$, and unitality is forced
by the prescribed image of $x^0$.
\end{proof}

\begin{remark}[What was not assumed]\label{rem:notassumed}
The proof never writes $\Phi(B)=\sum b_n\Phi(x)^n$. It applies $\Phi$ only
to the finite algebraic identity $B^2+x-1=0$. Consequently it does not
assume continuity, preservation of a valuation, preservation of order,
coefficientwise evaluation, injectivity, or preservation of infinite sums.
It also does not assume that arbitrary real coefficients are fixed. Only
unitality and the prescribed image of $x$ are used.
\end{remark}

\subsection{The smallest useful algebraic certificate}
The witness can be confined to a countable subring rather than the entire
power-series ring or a proper class.

\begin{proposition}\label{prop:small}
Let $A=\Q[x,B]\subseteq\Q[[x]]$. Then $A=\Q[B]$ is a countable polynomial
algebra in one transcendental element. Nevertheless, no unital homomorphism
$A\to F$ into an ordered field can send its distinguished element $x$ to a
number greater than $1$.
\end{proposition}
\begin{proof}
The relation $x=1-B^2$ gives $A=\Q[B]$. If $B$ were algebraic over $\Q$,
then $x=1-B^2$ would be algebraic over $\Q$, contrary to the transcendence of
the formal indeterminate $x$. Thus the map $\Q[Y]\to A$ sending $Y$ to $B$
is an isomorphism, under which $x$ corresponds to $1-Y^2$. The obstruction is
then just the inequality $1-y^2\le1$ in every ordered field. Finally, the set
of finite rational-coefficient polynomials is countable.
\end{proof}

This does not make $A$ an unembeddable ring: as a polynomial algebra it has
many embeddings into ordered fields. It is the requirement on the image of
one specified element, not the isomorphism type of the source, that fails.
Nor is $\Q[[x]]$ itself countable; only the smaller algebra $A$ is.

The infinite support of $B$ is needed only to \emph{place} the witness inside
the proposed domain. The final obstruction is a finite algebraic relation
between two elements of a countable ring. Consequently neither uncountable
supports nor delicate proper-class considerations enter the negative proof at
any point.

\subsection{Closed forms for the witness}
The coefficients also have the familiar expressions
\begin{equation}\label{eq:closed}
 b_0=1,\qquad
 b_n=(-1)^n\binom{1/2}{n}
 =-\frac{\binom{2n}{n}}{4^n(2n-1)}
 =-\frac{C_{n-1}}{2^{2n-1}}\quad(n\ge1),
\end{equation}
where $C_m=\frac{1}{m+1}\binom{2m}{m}$. To check the first expression directly,
let $D(x)$ denote the series with coefficients $(-1)^n\binom{1/2}{n}$.
Its coefficients obey
\[
 (n+1)[x^{n+1}]D=(n-\tfrac12)[x^n]D,
\]
so formal differentiation gives $(1-x)D'=-D/2$. For $E=D^2$ this yields
$(1-x)E'=-E$, with $E(0)=1$. If $E=\sum e_nx^n$, then
\[
 e_1=-1,\qquad (n+1)e_{n+1}=(n-1)e_n\quad(n\ge1).
\]
Thus $E=1-x$. The coefficient recursion for a square root with constant
coefficient $1$ is unique, so $D=B$. The remaining forms in \eqref{eq:closed}
follow by cancelling factorials. This argument is wholly formal; it does not
invoke a real analytic branch or a radius of convergence.

\paragraph{A second formal derivation.}
The same closed form can be verified by a quotient rather than a product, and
it is worth recording because the technique is exactly the one reused in
\cref{lem:H} below. Write $Q(x)=\sum_{n\ge0}q_nx^n$ with
$q_n=(-1)^n\binom{1/2}{n}$, so that the ratio of consecutive coefficients is
\[
 2(n+1)q_{n+1}=(2n-1)q_n,
 \qquad\text{equivalently}\qquad
 2(1-x)Q'(x)+Q(x)=0.
\]
The formal derivative of $Q(x)^2/(1-x)$ is then
\[
 \frac{2QQ'(1-x)+Q^2}{(1-x)^2}
 =\frac{Q\bigl(2(1-x)Q'+Q\bigr)}{(1-x)^2}=0,
\]
where $(1-x)^{-1}$ is the formal geometric series. In characteristic zero a
formal series with zero derivative is constant, and here the constant
coefficient is $1$. Hence $Q^2=1-x$, and the uniqueness clause of
\cref{lem:witness} identifies $Q$ with $B$. The two derivations differ in
substance and not only in bookkeeping: the first forms $E=D^2$ and reads off a
first-order recurrence for its coefficients, whereas the second divides by
$1-x$ and uses that a derivative-free series is constant.

\begin{table}[htbp]
\centering
\renewcommand{\arraystretch}{1.18}
\begin{tabular}{rcc}
\toprule
$n$ & $b_n=\coeff{x}{n}(1-x)^{1/2}$ & $h_n=\coeff{x}{n}(1-x)^{-1/2}$\\
\midrule
0 & $1$ & $1$\\
1 & $-1/2$ & $1/2$\\
2 & $-1/8$ & $3/8$\\
3 & $-1/16$ & $5/16$\\
4 & $-5/128$ & $35/128$\\
5 & $-7/256$ & $63/256$\\
6 & $-21/1024$ & $231/1024$\\
7 & $-33/2048$ & $429/2048$\\
8 & $-429/32768$ & $6435/32768$\\
\bottomrule
\end{tabular}
\caption{Exact rational coefficients of the two witness series. Both supports
have order type $\omega$. The second column is used in
\cref{sec:main,sec:inf}; the third is the display data for the
nonnegative-coefficient witness $H$ of \cref{sec:semiring}.}
\label{tab:coefficients}
\end{table}

\section{An obstruction with only nonnegative coefficients}\label{sec:semiring}
A natural reaction to \cref{sec:main} is to ask whether the negative
coefficients in $B$ are responsible for the contradiction. They are not.
Even retaining only nonnegative rational coefficients --- so that the source
has no additive inverses at all --- does not give a multiplicative domain
small enough to admit evaluation at $\omega$.

Let $\Qnn[[x]]$ denote the \emph{semiring} of formal power series with
nonnegative rational coefficients, under the ordinary coefficientwise addition
and the finite Cauchy product recalled at the start of
Section~\ref{sec:main}. Neither operation can leave $\Qnn[[x]]$, so this is a
well-defined commutative semiring, and it is not a ring: only the zero series
has an additive inverse. A unital semiring
homomorphism into an ordered field is required here to preserve only
$0$, $1$, $+$ and $\cdot$. It is not assumed to preserve any order, any
infinite sum, or anything else. In this section $G(x)=\sum_{n\ge0}x^n$ is the
formal geometric series, as in \eqref{eq:goodgeom} and \eqref{eq:badgeom}; the
letter $G$ without an argument continues to denote an exponent group, and the
letter $H$ with an argument denotes the witness series about to be defined,
never the Hahn field $\Hu{G}$.

\begin{lemma}[The nonnegative witness pair]\label{lem:H}
The series
\[
 H(x)=\sum_{n\ge0}h_nx^n,
 \qquad h_n=\frac{\binom{2n}{n}}{4^n},
 \qquad G(x)=\sum_{n\ge0}x^n
\]
both lie in $\Qnn[[x]]$ and satisfy
\begin{equation}\label{eq:positive-identities}
 H(x)^2=G(x),\qquad G(x)=1+xG(x).
\end{equation}
\end{lemma}
\begin{proof}
Every displayed coefficient is a positive rational, so both series lie in
$\Qnn[[x]]$. The second identity in \eqref{eq:positive-identities} is
immediate coefficientwise. For the first, the coefficient ratio
\[
 2(n+1)h_{n+1}=(2n+1)h_n
\]
is a direct computation with central binomial coefficients, and it says
exactly that $2(1-x)H'(x)=H(x)$. Calculating now in the containing
\emph{ring} $\Q[[x]]$, the formal derivative of $(1-x)H(x)^2$ is
\[
 -H^2+2(1-x)HH'=-H^2+H\cdot H=0.
\]
A formal series over a field of characteristic zero with zero derivative is
constant, and the constant coefficient of $(1-x)H^2$ is $1$. Hence
$(1-x)H(x)^2=1$. The element $1-x$ is a unit of $\Q[[x]]$ with inverse
$G(x)$, by the second identity, so $H(x)^2=G(x)$.

What has been proved is an equality of nonnegative-coefficient series, even
though subtraction was convenient in verifying it. The statement
\eqref{eq:positive-identities} mentions no negative number, and both of its
sides lie in $\Qnn[[x]]$.
\end{proof}

\begin{theorem}[Nonnegative-coefficient obstruction]\label{thm:semiring}
Let $K$ be an ordered field and let $u\in K$ satisfy $u\ge1$. There is no
unital semiring homomorphism
\[
 \psi:\Qnn[[x]]\longrightarrow K,\qquad \psi(x)=u.
\]
\end{theorem}
\begin{proof}
Suppose $\psi$ were such a map, and put $z=\psi(H)$ and $g=\psi(G)$. Applying
$\psi$ to the two identities of \eqref{eq:positive-identities} gives
\[
 g=z^2\ge0,\qquad g=1+ug.
\]
The second equation is an equation in the field $K$, where subtraction is
available even though it is unavailable in the source. Rearranging it gives
$(1-u)g=1$. Since $u\ge1$ and $g\ge0$, the left-hand side is nonpositive,
while the right-hand side is $1>0$. This is a contradiction.
\end{proof}

\begin{remark}[The two obstructions are incomparable]\label{rem:incomparable}
\Cref{thm:semiring} is not a corollary of \cref{thm:main}, nor conversely,
and neither makes the other redundant. It weakens the hypotheses of
\cref{thm:main} along two axes simultaneously: the source is only a semiring,
with no additive inverses and therefore with strictly fewer identities
available to the argument, and the excluded values are all $u\ge1$ rather than
only $u>1$. In the other direction, \cref{thm:main} applies to every unital
ring homomorphism defined on all of $\Q[[x]]$, a domain on which $\psi$ need
not be defined at all, and a ring homomorphism out of $\Q[[x]]$ does not
restrict to a semiring homomorphism claim about $\Qnn[[x]]$ unless one already
knows that the restriction lands where required. The reader should keep both.
\end{remark}

Only the subsemiring generated by $x$, $H$, $G$ and the nonnegative rational
constants is actually used. In particular, the images of arbitrary
positive-coefficient series were never assumed to be positive: the single
value that the proof needs to be nonnegative is $g$, and it is forced to be
nonnegative only because it is a square in the target. This is what keeps
\cref{thm:semiring} from being a disguised order-preservation hypothesis, and
it isolates the algebraic reason that a deliberately nonmonotone
interpretation still fails.

\section{The universal infinitesimal constraint}\label{sec:inf}
The obstruction is not peculiar to $\omega$ or to numbers greater than $1$.
The ring $\Q[[x]]$ forces a much stronger condition on the image of $x$.

\begin{definition}
In an ordered field $F$, an element $u$ is \emph{infinitesimal} if
\[
 |u|<\frac1n\qquad\text{for every integer }n\ge1.
\]
Here zero is included among the infinitesimals. A nonzero infinitesimal may
have either sign. When $F$ contains $\R$, this is equivalent to requiring
$|u|<r$ for every positive real $r$, since every positive real exceeds some
$1/n$ and every $1/n$ is a positive real.
\end{definition}

\begin{lemma}[Square roots of formal units]\label{lem:unitsquare}
For every $q\in\Q$, the series $1+qx$ is a square of a unit in $\Q[[x]]$.
\end{lemma}
\begin{proof}
Substitute $-qx$ for $x$ in $B(x)$: this simply multiplies its coefficient
$b_n$ by $(-q)^n$. The resulting formal series $B(-qx)$ satisfies
$B(-qx)^2=1+qx$. Its constant coefficient is $1$, so it is a unit. Indeed,
for a series $v=1+\sum_{n\ge1}v_nx^n$, the inverse coefficients are defined
by $w_0=1$ and $w_n=-\sum_{k=1}^{n}v_kw_{n-k}$.
\end{proof}

\Cref{lem:unitsquare} is all that \cref{thm:necessary} needs, and it has the
merit of using nothing beyond the single series $B$ already constructed. The
following companion statement is strictly more general in three directions at
once --- an arbitrary unit rather than $1+qx$, an arbitrary field of
characteristic zero rather than $\Q$, and uniqueness as well as existence ---
and it is what the arguments of \cref{prop:rational,thm:exact} use later.
Both are kept: the substitution proof of \cref{lem:unitsquare} is the one
reproduced in the accompanying one-page extraction, and it does not depend on
the recursion below.

\begin{lemma}[Square roots of units, general form]\label{lem:unit-root}
Let $k$ be a field of characteristic zero. Every $U\in1+xk[[x]]$ has a unique
square root in $1+xk[[x]]$. Moreover, every series in $k[[x]]$ with nonzero
constant coefficient is a unit.
\end{lemma}
\begin{proof}
Write $U=1+\sum_{n\ge1}u_nx^n$ and look for $V=1+\sum_{n\ge1}v_nx^n$ with
$V^2=U$. Comparing the coefficient of $x^n$ on the two sides gives
\[
 2v_n+\sum_{i=1}^{n-1}v_iv_{n-i}=u_n,
 \qquad\text{that is}\qquad
 v_n=\frac12\left(u_n-\sum_{i=1}^{n-1}v_iv_{n-i}\right).
\]
Division by $2$ is available because $\operatorname{char}k=0$. The
right-hand side involves only $v_1,\ldots,v_{n-1}$, so the equation determines
exactly one value of $v_n$ at each stage. This is simultaneously the existence
and the uniqueness assertion.

For the second claim, let $A=\sum_{n\ge0}a_nx^n$ with $a_0\ne0$ and seek
$D=\sum_{n\ge0}d_nx^n$ with $AD=1$. Put $d_0=a_0^{-1}$, and for $n\ge1$
require the coefficient of $x^n$ in $AD$ to vanish. In that equation the
coefficient multiplying the single new unknown $d_n$ is $a_0$, which is
invertible, so again exactly one value is determined at each stage.
\end{proof}

\begin{theorem}[Necessary infinitesimal condition]\label{thm:necessary}
Every unital ring homomorphism $\phi:\Q[[x]]\to F$ into an ordered field
satisfies
\begin{equation}\label{eq:infbound}
 -\frac1n<\phi(x)<\frac1n\qquad(n\ge1).
\end{equation}
In particular, no such homomorphism sends $x$ to a nonzero real number or
to a non-infinitesimal element of a non-Archimedean ordered field.
\end{theorem}
\begin{proof}
For each $n\ge1$, both $1+nx$ and $1-nx$ are squares of units by
\cref{lem:unitsquare}. A unital homomorphism maps a unit to a nonzero element:
if $vw=1$, then $\phi(v)\phi(w)=1$. The images of the two unit squares are
therefore strictly positive. With $u=\phi(x)$ this gives
\[
 1+nu>0,\qquad 1-nu>0,
\]
which is exactly \eqref{eq:infbound}.
\end{proof}

\begin{proposition}[The kernel dichotomy]\label{prop:kernel}
Let $\phi:\Q[[x]]\to F$ be a unital homomorphism into a field. If
$\phi(x)=0$, then $\phi$ is the constant-term map. If $\phi(x)\ne0$,
then $\phi$ is injective.
\end{proposition}
\begin{proof}
Unitality forces the rational constants to be fixed. If $\phi(x)=0$, write
$f=f_0+xg$ to obtain $\phi(f)=f_0$. If $\phi(x)\ne0$, every nonzero series
has a factorization $f=x^m v$ with $m\ge0$ and $v$ a unit. Both factors
have nonzero images, so $\phi(f)\ne0$.
\end{proof}

Consequently, if the target is Archimedean, the only unital homomorphism
$\Q[[x]]\to F$ is the constant-term map. In particular, every $\R$-algebra
homomorphism $\R[[x]]\to\R$ is constant-coefficient evaluation: the value of
$x$ must be an infinitesimal real number, hence zero, and then the
decomposition $f=f(0)+xg$ shows that the value of $x$ determines the whole
map. For the Laurent-series field $\Q((x))$ even the constant-term map is
unavailable: $x$ is invertible, so its image must be a nonzero infinitesimal.
These deductions use only the field axioms and the preceding lemmas.

The corresponding question for a non-Archimedean target needs more care, and
the rest of this section is devoted to it. There, nonzero infinitesimal values
of $x$ genuinely do occur, so \cref{prop:kernel} leaves a real alternative
open; and the infinite sums that describe the resulting maps must be read
formally, through the support calculus of \cref{app:summability}, rather than
as limits in a topology on the target. We write $\seqsum$ for a sum defined in
that way whenever the distinction matters, and \cref{sec:notlimit} shows by an
explicit computation that the distinction is not cosmetic.

\subsection{Exact converse for surreal targets}
For $\No$ the necessary condition is also sufficient. The relevant notion
of infinite sum is Hahn summability, not convergence in an unspecified topology.
A family of surreal normal forms is summable when the union of their exponent
supports is reverse well-ordered and each exponent occurs in only finitely
many members. Such a sum is computed coefficientwise
\citep[Definition 2.9]{BerarducciMantova2018}.

\begin{theorem}[Exactly which values of $x$ occur]\label{thm:exact}
For a surreal number $u$, the following are equivalent:
\begin{enumerate}[label=(\roman*)]
\item There exists a unital homomorphism $\Q[[x]]\to\No$ sending $x$ to $u$.
\item $u$ is infinitesimal, including the possibility $u=0$.
\item There exists a coefficient-fixing, strongly additive homomorphism
$\R[[x]]\to\No$ sending $x$ to $u$.
\item There exists an $\R$-algebra homomorphism $\R[[x]]\to\No$ sending $x$
to $u$.
\end{enumerate}
For nonzero $u$, a homomorphism in (iii) is given by the Hahn sum
\begin{equation}\label{eq:ev}
 \ev_u\left(\sum_{n\ge0}c_nx^n\right)=\seqsum_{n\ge0}c_nu^n.
\end{equation}
It is injective. For $u=0$, the constant-term map gives the required map.
\end{theorem}
\begin{proof}
The implication (i)$\Rightarrow$(ii) is \cref{thm:necessary}, and
(iii)$\Rightarrow$(i) follows by restriction. Clause~(iv) sits between the
other two: a coefficient-fixing strongly additive homomorphism is in
particular $\R$-linear, so (iii)$\Rightarrow$(iv); and an $\R$-algebra
homomorphism is unital, so its restriction to $\Q[[x]]$ witnesses~(i) and
(iv)$\Rightarrow$(i). No strong additivity is assumed in~(iv). Suppose (ii) holds and $u\ne0$.
All exponents in the normal form of $u$ are strictly negative. Negating this
support gives a well-ordered set $A\subseteq\No_{>0}$.

The positive-support summability lemma proved in
Appendix~\ref{app:summability} says that the additive monoid generated by
$A$ is well-ordered and each exponent has only finitely many representations
as a finite word over $A$. Reflecting exponents back, this shows that the
family $(c_nu^n)_{n\ge0}$ is summable, regardless of the sizes of the real
coefficients $c_n$.

The same finite-representation statement justifies expanding products and
regrouping by total degree. Thus, if $f=\sum c_nx^n$ and $g=\sum d_nx^n$,
\[
 \ev_u(f)\ev_u(g)
 =\sum_{i,j\ge0}c_id_ju^{i+j}
 =\sum_{n\ge0}\left(\sum_{i+j=n}c_id_j\right)u^n
 =\ev_u(fg).
\]
Each coefficient of this regrouping receives only finitely many contributions:
a finite representing word has only finitely many places at which it can be
split into two words. Addition and preservation of constants are immediate.

For strong additivity, take a source-summable family of power series. At each
degree only finitely many members contribute. At a fixed target exponent only
finitely many degrees can contribute, by the same positive-support lemma.
The two finiteness conditions allow the sums to be interchanged, and their
union support lies in a fixed reverse well-ordered set. Hence \eqref{eq:ev}
preserves every source-summable family.

Finally, for $f\ne0$ write $f=x^m v$, where $v$ is a unit in $\R[[x]]$.
The map just constructed sends $u^m$ and $v$ to nonzero elements, so it is
injective. The argument for $u=0$ is the direct constant-term calculation.
\end{proof}

\begin{remark}[Injectivity again, by leading terms]\label{rem:leading-inj}
Injectivity of $\ev_u$ also follows from a direct calculation with normal
forms, which is worth recording because it produces the leading term of the
image rather than merely certifying that it is nonzero, and because it uses
no factorization of $f$ and hence neither \cref{lem:unit-root} nor the unit
structure of $\R[[x]]$. Let $a<0$ be the largest exponent occurring in the
normal form of the nonzero infinitesimal $u$, and let $r\ne0$ be the
corresponding leading real coefficient. Let $f=\sum_{n\ge0}c_nx^n$ be nonzero
and let $m$ be the least index with $c_m\ne0$. The term $c_mu^m$ then has
leading exponent $ma$ and leading coefficient $c_mr^m\ne0$, by
\eqref{eq:lead}. For every $n>m$, each exponent occurring in $u^n$ is at most
$na$, and $na<ma$ because $a<0$ and $n>m$. So no later term of the Hahn sum
\eqref{eq:ev} reaches the exponent $ma$, and the leading term
$c_mr^m\omega^{ma}$ cannot cancel. Hence $\ev_u(f)\ne0$.
\end{remark}

Strong additivity makes \eqref{eq:ev} the unique coefficient-fixing strongly
additive map with that value of $x$, because every power series is a summable
family of its monomials. We do \emph{not} infer uniqueness among arbitrary
abstract ring homomorphisms. Describing the maps that fail to respect infinite
sums is a separate question, and the distinction between arbitrary and
strongly additive maps is a substantive one in the surreal setting; compare
\citep[Sections 3 and 4]{KKS}. The successful termwise construction is the
familiar change-of-value-group construction for Hahn fields, written in the
opposite exponent convention; compare \citep[Section 2.3]{KKS}. In particular,
the nonexistence theorem
requires far less structure than the positive construction uses.

There is no contradiction with classical analytic convergence. A series may
have zero radius of convergence over the real numbers and still admit
\eqref{eq:ev}: its real coefficients change the values at exponents, but they
do not change the well-ordering or finite-occurrence condition on supports.

\subsection{Summability is not convergence in the full order topology}\label{sec:notlimit}
The sums above are defined by the support calculus of
\cref{app:summability}, not as limits. That is not a matter of taste. Consider
the entirely legitimate Hahn sum
\[
 s=\seqsum_{n\ge0}\omega^{-n}=\frac{1}{1-\omega^{-1}},
\]
which is \eqref{eq:goodgeom} again, and let $s_N=\sum_{n=0}^{N}\omega^{-n}$ be
its $N$-th partial sum, a finite and therefore unambiguous surreal number.
Summing the remaining geometric tail gives
\[
 s-s_N=\seqsum_{n\ge N+1}\omega^{-n}
      =\frac{\omega^{-(N+1)}}{1-\omega^{-1}}
      >\omega^{-\omega}\qquad(N\ge0),
\]
the last inequality because $0<1-\omega^{-1}<1$ forces
$(1-\omega^{-1})^{-1}>1$, while $-(N+1)>-\omega$ for every natural number $N$
and the omega-map is strictly increasing by \eqref{eq:omega}. The single fixed positive
surreal $\omega^{-\omega}$ therefore determines a neighbourhood of $s$ in the
order topology that \emph{no} partial sum $s_N$ ever enters. The family is
summable and its sum is $s$, yet the sequence of partial sums does not
converge to $s$ in any order-theoretic sense. This is why
\cref{thm:exact} is stated in terms of Hahn summability and why no appeal to
convergence is made anywhere in this note.

\section{The order hidden in the full Hahn field}\label{sec:order}
The main proof deliberately avoids assuming source-order preservation. In the
full Hahn field, however, order preservation is not an optional extra: it can
be recovered from squares.

\subsection{Leading terms and the canonical order}
For $s\ne0$ in $\Hu{G}$ let
\[
 \vexp(s)=\min\supp(s),\qquad \lc(s)=s_{\vexp(s)}.
\]
Declare $s>0$ if $\lc(s)>0$. For two nonzero series,
\begin{equation}\label{eq:lead}
 \vexp(st)=\vexp(s)+\vexp(t),\qquad
 \lc(st)=\lc(s)\lc(t).
\end{equation}
Indeed, the smallest possible exponent of a product is the sum of the two
smallest exponents; it is attained only by that pair. The coefficient cannot
cancel. The sum of two positive series is positive as well: at the smaller
leading exponent its coefficient is positive, or, when the exponents agree,
it is the sum of two positive coefficients. Thus this is a compatible total
ordering.

For $a>0$, the monomial $x^a$ is positive, and for every positive real $r$
the series $r-x^a$ has leading coefficient $r$ at exponent $0$. Therefore
\begin{equation}\label{eq:source-small}
 0<x^a<r\qquad(a>0,\ r\in\R_{>0}).
\end{equation}
Larger exponents encode smaller magnitudes in this source convention.

\begin{lemma}[Every canonically positive element is a square]\label{lem:positivesquare}
Suppose $G$ is $2$-divisible: every $a\in G$ has a half $a/2\in G$.
Then the positive elements of $\Hu{G}$ are exactly its nonzero squares.
\end{lemma}
\begin{proof}
A nonzero square has positive leading coefficient by \eqref{eq:lead}.
Conversely, if $s>0$, factor its leading term:
\[
 s=r x^a(1+\epsilon),\qquad r>0,
 \qquad \supp(\epsilon)\subseteq G_{>0}.
\]
By the summability lemma in Appendix~\ref{app:summability},
\[
 q=\sqrt r\,x^{a/2}\sum_{n\ge0}\binom{1/2}{n}\epsilon^n
\]
is a well-defined Hahn series. Formal substitution of $\epsilon$ into the
power-series square-root identity is legitimate coefficientwise, again by
finite representability. It gives $q^2=s$.
\end{proof}

\begin{theorem}[Order is algebraically forced]\label{thm:orderforced}
For $2$-divisible $G$, the field $\Hu{G}$ has only one ordering compatible
with its field operations. Every unital ring homomorphism from $\Hu{G}$
into an ordered field preserves that ordering.
\end{theorem}
\begin{proof}
In any ordered field every nonzero square is positive. By
\cref{lem:positivesquare}, each nonzero element is either a square or the
negative of a square according to the canonical order. Hence every compatible
ordering has the same positive cone.

A unital homomorphism out of a field is injective: if $s\ne0$, then
$\phi(s)\phi(s^{-1})=1$. A positive source difference is a nonzero square,
so its image is a nonzero square and is positive in the target. This is
exactly order preservation.
\end{proof}

The fact that $\Hu{G}$ is a field follows from the formal geometric inverse.
Concretely, writing a nonzero $s$ as $s=rx^a(1+\epsilon)$ with $r\ne0$,
$a=\vexp(s)$ and $\supp(\epsilon)\subseteq G_{>0}$, the inverse is
\[
 s^{-1}=r^{-1}x^{-a}\sum_{n\ge0}(-\epsilon)^n,
\]
a sum that is legitimate by the support calculus rather than by convergence.
The full proof, valid for an arbitrary linearly ordered abelian group $G$, is
\cref{prop:field} in Appendix~\ref{app:summability}; it is stated there in that
generality rather than only for $G=\No$. No real-closedness
theorem is needed for \cref{thm:orderforced}, and having all positive elements
square should not, by itself, be confused with real closedness.

Applied to $G=\No$, this gives a second explanation of the impossibility:
$x$ is a positive infinitesimal in the source, but the requested image is
positive infinite. One cannot rescue the prescription by asking for a
non-monotone ring homomorphism, since there is no such freedom on this field.
The short proof in Section~\ref{sec:main} is nevertheless stronger as a local
obstruction: it works before adjoining any fractional exponents and without
proving uniqueness of the source ordering.

\section{The classical repair: reverse the exponents}\label{sec:repair}
The source field is not intrinsically incompatible with the surreals. The
correct normal-form identification reverses the exponent order. This is a
classical construction: van den Dries and Ehrlich explicitly display its
restriction to any set-sized surreal exponent group on p.~176 of
\citet{DriesEhrlich2001}.

\begin{theorem}[Exponent-reversing isomorphism]\label{thm:reversal}
The map
\begin{equation}\label{eq:J}
 J:\Hfull\longrightarrow\No,
 \qquad J\left(\sum_{a\in A}r_ax^a\right)
       =\sum_{a\in A}r_a\omega^{-a}
\end{equation}
is a coefficient-fixing ordered-field isomorphism. It is strongly additive
and satisfies $J(x)=\omega^{-1}$.
\end{theorem}
\begin{proof}
If $A$ is well-ordered, $-A$ is reverse well-ordered. The expression on the
right of \eqref{eq:J} is therefore a legitimate Conway normal form and so
defines a unique surreal number. Negation of exponents is injective, hence
there is no identification or cancellation between distinct source monomials.

The map preserves addition coefficientwise. For multiplication, the
coefficient at target exponent $-c$ is the finite sum over pairs satisfying
$-a-b=-c$, equivalently $a+b=c$. This is exactly the source convolution.
It preserves $1$ and fixes every real coefficient.

A nonzero source leading term $r x^a$ becomes the target leading term
$r\omega^{-a}$, with the same coefficient, so signs are preserved. Finally,
any surreal normal form $\sum_b c_b\omega^b$ has the inverse image
$\sum_b c_bx^{-b}$; the reflected support is well-ordered. This proves
surjectivity as well as injectivity. For a summable family, reflection
preserves the finite-occurrence property and changes the union support from
well-ordered to reverse well-ordered. Thus $J$ is strongly additive.
\end{proof}

\subsection{Two explicit evaluations}
For the formal geometric series $G(x)=\sum_{n\ge0}x^n$ we have
\begin{equation}\label{eq:goodgeom}
 J(G)=\sum_{n\ge0}\omega^{-n}=\frac{1}{1-\omega^{-1}}.
\end{equation}
For the obstructing square-root series,
\begin{equation}\label{eq:goodroot}
 J(B)=1-\frac{\omega^{-1}}2-\frac{\omega^{-2}}8
       -\frac{\omega^{-3}}{16}-\cdots,
 \qquad J(B)^2=1-\omega^{-1}.
\end{equation}
The right side of the square identity is positive. Its displayed root has
positive leading coefficient and is therefore the positive square root.
There is no order conflict.

\subsection{Exponent reversal is not field-order reversal}
The word ``reversing'' in \cref{thm:reversal} applies to the exponent parameter,
not to the ordering of field elements. The map $a\mapsto-a$ reverses the
exponent order, while the map $J$ preserves field order. This is exactly what
is needed: source positive exponents mean small monomials, whereas positive
exponents in $\omega^a$ mean large monomials.

Replacing $x$ by $\omega^{-1}$ in this way is a repair of the prescription,
not a positive solution of Problem 7.7. Its demanded images are $\omega^a$,
not $\omega^{-a}$.

\section{Classification of pure monomial prescriptions}\label{sec:classification}
We can classify the monomial data compatible with a homomorphism without
assuming in advance that the homomorphism is strongly additive.

\begin{definition}
A map between Hahn-series structures is \emph{strongly additive} if it takes
every summable family to a summable family and commutes with its sum.
For $\Hu{G}$, summability means that the union of exponent supports is
well-ordered and each exponent occurs in only finitely many family members.
For $\No$, the union support is instead reverse well-ordered.
\end{definition}

\begin{theorem}[Complete existence criterion for pure monomial data]\label{thm:classification}
Let $G$ be a set-sized linearly ordered abelian group, or let $G=\No$ with
set-sized supports. Let $f:G\to\No$ be a function. The following conditions
are equivalent:
\begin{enumerate}[label=(\roman*)]
\item There exists a unital ring homomorphism $\Phi:\Hu{G}\to\No$ such that
$\Phi(x^a)=\omega^{f(a)}$ for every $a\in G$.
\item The map $f$ is additive and strictly order-reversing:
\begin{equation}\label{eq:anti}
 f(a+b)=f(a)+f(b),\qquad a<b\Longrightarrow f(a)>f(b).
\end{equation}
\item There exists a coefficient-fixing, strongly additive homomorphism with
those monomial images.
\end{enumerate}
When these conditions hold, the map in (iii) is unique and is the embedding
\begin{equation}\label{eq:Tf}
 T_f\left(\sum_{a\in A}r_ax^a\right)
       =\sum_{a\in A}r_a\omega^{f(a)}.
\end{equation}
Moreover, $T_f$ is onto $\No$ if and only if $f$ is onto $\No$.
\end{theorem}
\begin{proof}
\emph{Necessity of additivity.}
Unitality and the identity $x^{a+b}=x^ax^b$ give
\[
 \omega^{f(0)}=1=\omega^0,
 \qquad
 \omega^{f(a+b)}=\omega^{f(a)}\omega^{f(b)}
               =\omega^{f(a)+f(b)}.
\]
Injectivity of the omega-map implies $f(0)=0$ and additivity.

\emph{Necessity of order reversal.}
Fix $a>0$. Substitution of the monomial $x^a$ for a new formal variable $z$
embeds $\Q[[z]]$ into $\Hu{G}$:
\[
 \sum_{n\ge0}q_nz^n\longmapsto\sum_{n\ge0}q_nx^{na}.
\]
The support is well-ordered because $0,a,2a,\ldots$ is strictly increasing;
convolution agrees with ordinary power-series multiplication. Compose this
embedding with $\Phi$. By \cref{thm:necessary}, its image of $z$, namely
$\omega^{f(a)}$, is infinitesimal. A pure Conway monomial is a positive
infinitesimal exactly when its exponent is negative. Thus $f(a)<0$ for
all $a>0$. If $a<b$, additivity gives
\[
 f(b)-f(a)=f(b-a)<0,
\]
which is strict order reversal. This argument does not require $G$ to be
divisible or $\Phi$ to fix the real coefficients.

\emph{Sufficiency and construction.}
Assume (ii). The image $f(A)$ of every well-ordered support is reverse
well-ordered, and $f$ is injective. Therefore \eqref{eq:Tf} is a valid
normal form. Addition is preserved coefficientwise. For products, additivity
identifies $f(a)+f(b)$ with $f(a+b)$, and injectivity ensures that equality
of image exponents is equivalent to equality of their source sums. Hence
the finite convolution coefficients agree.

The construction fixes the real constants and preserves $1$. A nonzero
coefficient remains nonzero at a unique image exponent, so $T_f$ is injective.
A summable family has a well-ordered union support and finite occurrence at
each source exponent. Its image has a reverse well-ordered union support
and the identical occurrence counts. This proves strong additivity and
(iii). The implication (iii)$\Rightarrow$(i) is immediate.

\emph{Uniqueness under strong additivity.}
Every source series is the sum of its summable family of coefficient-monomial
terms. A coefficient-fixing strongly additive map must therefore have exactly
the value in \eqref{eq:Tf}.

\emph{Surjectivity.}
If $f$ is onto, apply $f^{-1}$ to the reverse well-ordered support of an
arbitrary surreal normal form; the resulting support is well-ordered and
provides a preimage. Conversely, if $T_f$ is onto, every monomial $\omega^b$
is in its image. Normal-form uniqueness in \eqref{eq:Tf} forces
$b=f(a)$ for a source exponent $a$.
\end{proof}

The equivalence of (i) and (ii) is an existence classification even for
ordinary, potentially non-strong homomorphisms. The uniqueness assertion,
by contrast, applies only to the coefficient-fixing strongly additive maps.
These are distinct statements and should not be conflated.

\subsection{Examples and boundary cases}
\begin{example}[A family of isomorphisms]
For any fixed $c\in\No_{>0}$, set $f_c(a)=-ca$, where multiplication is
surreal field multiplication. This map is additive, strictly order-reversing,
and bijective on $\No$. Thus
\begin{equation}\label{eq:Tc}
 T_c\left(\sum_ar_ax^a\right)=\sum_ar_a\omega^{-ca}
\end{equation}
is a coefficient-fixing ordered-field isomorphism $\Hfull\to\No$.
Its inverse is obtained by replacing an exponent $b$ by $-b/c$. In particular,
$x$ goes to $\omega^{-c}$, a positive infinitesimal even when $c$ itself is
infinitesimal.
\end{example}

\begin{example}[The forbidden identity exponent map]
The prescription in Problem 7.7 is $f(a)=a$. It is additive but order-preserving,
not order-reversing. It fails the criterion on the single positive exponent
$a=1$. Thus the general classification recovers the short contradiction.
\end{example}

\begin{example}[Zero is not a Hahn-field monomial image]
In \cref{thm:exact}, sending $x$ to zero is allowed because $x$ is not a unit
of $\Q[[x]]$. In $\Hu{G}$ every monomial is a unit, with inverse $x^{-a}$.
It cannot have image zero under a unital homomorphism. This accounts for an
important difference between the power-series ring and a full Hahn field.
\end{example}

An increasing exponent sequence can have arbitrary ordinal length, but an
order-reversing $f$ handles all such supports simultaneously. No bound on
their ordinal types or on the birthday of their surreal exponents is needed.
That is the order-theoretic content of the sufficient direction.

\section{Why rational evaluation at infinity does not contradict the result}\label{sec:partial}
The failure is not that every formal expression becomes meaningless at
$\omega$. Large useful subfields admit precisely that evaluation. What fails
is extending it to the entire power-series or Hahn field.

\subsection{The rational subfield}
\begin{proposition}\label{prop:rational}
The assignment $x\mapsto\omega$ extends to an injective, real-coefficient-fixing
homomorphism
\[
 \ev_\omega:\R(x)\longrightarrow\No,
 \qquad \frac{p(x)}{q(x)}\longmapsto\frac{p(\omega)}{q(\omega)}.
\]
Here $\R(x)$ is regarded as a subfield of $\R((x))$ by its Laurent expansion
at zero.
\end{proposition}
\begin{proof}
First, the Laurent embedding is genuinely available and not merely asserted.
Let $q(x)$ be a nonzero polynomial and factor it as $q(x)=x^mq_0(x)$ with
$q_0(0)\ne0$. By the second clause of \cref{lem:unit-root}, $q_0$ is a unit of
$\R[[x]]$, so $1/q=x^{-m}q_0^{-1}$ lies in $\R((x))$. Hence every rational
function has a formal Laurent expansion at zero, and $\R(x)\subseteq\R((x))$
as claimed.

For the evaluation, let $p(X)=a_dX^d+\cdots+a_0$ be a nonzero polynomial with
$a_d\ne0$. The surreal number $p(\omega)$ has leading normal-form term
$a_d\omega^d$, which dominates the finitely many lower monomials in absolute
value. It is therefore nonzero. Thus $\omega$ is transcendental over the real
coefficient field, and substitution extends from the polynomial ring to its
fraction field. Well-definedness follows by cross multiplication and
injectivity follows from the same nonvanishing statement.
\end{proof}

The element $G(x)=1+x+x^2+\cdots$ belongs to this rational subfield because
\[
 (1-x)G(x)=1.
\]
Hence the rational evaluation does consistently assign it
\begin{equation}\label{eq:badgeom}
 \ev_\omega(G)=\frac{1}{1-\omega}
   =-\sum_{n\ge1}\omega^{-n}<0.
\end{equation}
The last series is a legitimate decreasing-exponent normal form.
In particular, the geometric-series calculation by itself is not a
contradiction to the existence of an ordinary ring homomorphism. Positivity
of that source series need not be preserved on the rational subfield. That
field carries more than one compatible field ordering: there is one in which
$x$ is a positive infinitesimal, namely the one induced from $\R((x))$ by
leading terms, and another in which $x$ is positive and infinite, namely the
one pulled back along $x\mapsto\omega$. An algebraic embedding is free to
switch between two such orderings, and $\ev_\omega$ does exactly that. What
cannot be transported by any embedding is a \emph{square root} of a source
element whose image is negative, since a square stays a square.

The new obstruction appears upon adjoining $B$: the compatible image would
have to solve $Y^2=1-\omega$. Thus \cref{prop:rational} cannot extend to
$\R(x,B)$, although it is perfectly valid on $\R(x)$.

\subsection{A contrasting quadratic extension}
Let $C(x)=B(-x)\in\R[[x]]$, so that $C(x)^2=1+x$ and $C(0)=1$.
Unlike $B$, this series can be adjoined to the rational domain of evaluation.
The polynomial $Y^2-(1+x)$ is irreducible over $\R(x)$: the rational function
$1+x$ has a zero of odd multiplicity at $x=-1$, whereas a square of a
rational function has even multiplicity at every zero and pole.

Since $1+\omega>0$, it has a positive square root in $\No$. Concretely,
\[
 \sqrt{1+\omega}
 =\omega^{1/2}\sum_{n\ge0}\binom{1/2}{n}\omega^{-n},
\]
which is a valid normal form with square $1+\omega$. Thus
$x\mapsto\omega$ and $C\mapsto\sqrt{1+\omega}$ define an embedding of
$\R(x,C)$ into $\No$.

This embedding does not evaluate the original Taylor series of $C$ term by
term at $\omega$. Instead it respects the finite algebraic relation and
uses a different, decreasing-exponent expansion of its chosen value.
Moreover, either sign of the square root gives a field embedding, illustrating
why algebraic extensions may involve choices even when $x$ is fixed.

\begin{proposition}[Criterion for a simple algebraic extension]\label{prop:algebraic}
Let $p(Y)\in\R(x)[Y]$ be monic and irreducible, and let
$L=\R(x)[Y]/(p)$. The evaluation $x\mapsto\omega$ extends to an embedding
$L\to\No$ if and only if the polynomial $p^\omega(Y)$ obtained by evaluating
its rational-function coefficients at $\omega$ has a root in $\No$.
\end{proposition}
\begin{proof}
Any extension sends the residue class of $Y$ to such a root. Conversely, a
root induces a unital homomorphism from $\R(x)[Y]$ to $\No$ that annihilates
$(p)$. It factors through $L$, which is a field because $p$ is irreducible.
A unital homomorphism from a field has zero kernel, so the induced map is
an embedding.
\end{proof}

The pair $p=Y^2-(1-x)$ and $p=Y^2-(1+x)$ therefore separates two very similar
formal Taylor series: the first algebraic extension is forbidden at positive
infinity, while the second is permitted. This is a field-theoretic obstruction,
not merely a rule about rearranging divergent series.

It should not be inferred that $\R(x)$ is a \emph{largest} subfield on which
evaluation at $\omega$ is admissible, and no such claim is made here. The
preceding paragraphs exhibit a proper algebraic extension $\R(x,C)$ of it on
which the evaluation still exists, precisely because $1+\omega$ is positive
and so has a square root in $\No$. What \cref{thm:main} identifies is one
particular forbidden algebraic extension, not every algebraic extension; we do
not classify the maximal admissible subfields.

\section{Verification, attribution, and the boundary of the result}\label{sec:audit}

\subsection{What has been proved}
Theorem~\ref{thm:main} gives a complete negative answer to the precise ring
homomorphism question in Problem 7.7. Its proof is an explicit algebraic
contradiction, not a computational search with no examples found. It does not
settle the separate questions in the surrounding paper about other possible
infinitary operations on surreal numbers or combinatorial games.

The additional conclusions are exact statements with distinct scopes:
Theorem~\ref{thm:semiring} strengthens the obstruction to a semiring source
with no additive inverses and to all values $u\ge1$;
Theorem~\ref{thm:exact} classifies possible values of one power-series variable
in $\No$; Theorem~\ref{thm:classification} classifies pure monomial data for
Hahn-field homomorphisms; Theorem~\ref{thm:reversal} recalls the standard
exponent-reversing identification. We do not classify all arbitrary embeddings
of the surreal field, all non-monomial Hahn substitutions, or all maximal
subfields on which evaluation at $\omega$ is possible.

Stated as a ledger, exactly four things are proved: the nonexistence theorem
for the posed prescription; its nonnegative-coefficient semiring
strengthening; the image criterion for a single power-series variable; and the
existence criterion for prescribed monomial images. Everything else in this
note is either classical background, a worked example, or a scope
disclaimer.

\subsection{The normalization cannot be omitted}
If one merely asks for \emph{some} unital ring map $\Hfull\to\No$, one exists:
\cref{thm:reversal} supplies an isomorphism. The condition
$\Phi(x^a)=\omega^a$ is therefore essential to the problem and not a piece of
decoration on it. In the other direction, the condition is used only very
weakly: the much smaller pair of requirements $\Phi(1)=1$ and
$\Phi(x)=\omega$ already contradicts the subring certificate of
\cref{prop:small}, so the values at all other monomials are never consulted.
Finally, dropping unitality does admit a map, namely the zero map, but that
map fails the monomial condition already at $a=0$, where it gives $0$ rather
than $\omega^0=1$. There is no gap between these three observations for a
counterexample to hide in.

\subsection{Finite supports, polynomials, and truncations}
The ring of \emph{finite} formal sums of monomials can be evaluated termwise
at $\omega$, since no infinite support needs interpretation; $\R[x]$ and the
rational domain of \cref{prop:rational} present no obstruction whatever. It is
instructive that $B$ is absent from the polynomial ring for a reason that has
nothing to do with the target: a square polynomial has even degree, whereas
$1-x$ has degree one. The question thus fails only once the domain is enlarged
far enough to contain the specific algebraic power series $B$, and not merely
because formal monomials with infinite supports are permitted.

Nor can the contradiction be reached, or evaded, by evaluating truncations.
For $B_N=\sum_{n\le N}b_nx^n$ the identity $B_N^2=1-x$ is simply
\emph{false}; it holds only modulo $x^{N+1}$. At $x=\omega$ the discarded
high-degree terms are not small --- they are the dominant ones. The
contradiction in \cref{thm:main} is about the existence of an exact
homomorphism on the series ring, not about a limiting calculation with
polynomials, and a reader who reconstructs it as the latter will find it
unconvincing for the right reason.

\subsection{All support lengths are set-sized}
Although the exponents range over a proper class, each individual series in
\eqref{eq:source} has ordinal-indexed and therefore set-sized support. No sum
indexed by the entire class of ordinals is introduced anywhere. The
positive-support arguments of Appendix~\ref{app:summability} use only sets of
exponents and finite words over those sets, and the obstruction itself lives
entirely inside the set-sized ring $\Q[[x]]$ together with a hypothetical
image in an ordered field. Class-sized notation, where it appears, can be
handled in a standard class theory such as NBG. Consequently the negative
result cannot be avoided by changing the treatment of proper classes.

\subsection{The target must be formally real for this obstruction}
The argument uses exactly one property of $\No$ beyond its ring structure:
that a square is nonnegative in an ordered field. It does not claim
nonexistence for every imaginable target ring. If a target contains an element
whose square is $1-u$, this particular certificate produces no contradiction
there. Allowing non-real values would of course change the problem, which is
posed with a surreal-valued map; no assertion about a global homomorphism into
a complexified target is made or implied here.

\subsection{Classical ingredients versus the application}
The formal binomial square root, the leading-term order on Hahn fields,
positive-support summability, the omega-map, and exponent reversal are
classical. In particular, \eqref{eq:J} is not being proposed as a new
isomorphism; its set-sized versions occur explicitly in
\citet[p.~176]{DriesEhrlich2001}. The word-embedding result used in the appendix
is Higman's classical theorem \citep{Higman1952}; we give the special-case proof
needed here.

The contribution of this note is the direct application of the square-root
identity to the specifically posed evaluation question and the organized
account of its exact consequences. The elementary nature of the obstruction
is a reason to make its assumptions unusually explicit, not a reason to
claim a new general theory of Hahn fields.

\subsection{Source version and priority qualification}
The inspected arXiv record lists the revised version dated 30 April 2026 as
its latest version. The question is still explicit there as Problem 7.7;
the earlier formulation is Problem 6.8 in the May 2025 version. Numbering in
this note refers to the extended preprint, not to an assumption that the
question occurs in the shortened journal version \citep{Lipparini2026}.

A focused public-source search was carried out on 20 September 2026 for
subsequent resolutions and for the specific Hahn/surreal homomorphism
formulation. No prior resolution was located in that search. This is not an
exhaustive priority certification, and it cannot rule out an unpublished
observation or an unindexed later discussion. The proof stands independently
of that bibliographic uncertainty. The manuscript has not been independently
refereed or checked in a proof assistant.

\subsection{Exact finite regression checks}
The archive includes \texttt{code/verify.py}, using only Python's standard
library and exact \texttt{fractions.Fraction} arithmetic. No floating-point
number is created anywhere in it. The executed default run verified the square
identity through degree $256$, that is,
\[
 B_{256}(x)^2\equiv1-x\pmod{x^{257}},
\]
and the two semiring identities of \cref{lem:H} to the same modulus,
\[
 H(x)^2\equiv G(x),\qquad B(x)H(x)\equiv1,\qquad (1-x)G(x)\equiv1
 \pmod{x^{257}}.
\]
It compared $257$ coefficients of $B$ obtained independently from five routes:
the quadratic recursion of \cref{lem:witness}, the general unit square-root
recursion of \cref{lem:unit-root} applied to $1-x$, the first-order binomial
recursion, the Catalan closed formula, and Catalan numbers regenerated from
the convolution recurrence $C_m=\sum_iC_iC_{m-1-i}$. It verified that all
computed nonconstant coefficients of $B$ are negative and dyadic, that all
$257$ computed coefficients of $H$ are strictly positive, and that both
differential recurrences $2(n+1)b_{n+1}=(2n-1)b_n$ and
$2(n+1)h_{n+1}=(2n+1)h_n$ hold at every checked degree.

For unit square roots there are two independent batteries at degree $64$: the
$32$ scaled checks for $1\pm nx$ with $1\le n\le16$, obtained by scaling the
coefficients of $B$, and $12$ checks for $1+cx$ at the signed rational values
$c\in\{\pm1,\pm2,\pm3,\pm\tfrac12,\pm\tfrac23,\pm\tfrac75\}$, obtained from
the general routine. The rational values exercise code that the integer
scalings do not.

For exponent reversal there are likewise two harnesses, and both are run.
The first performs $1000$ deterministic randomized tests of finite sparse
rational-exponent polynomials at randomly drawn positive rational scales, with
seed $20260920$. The second performs $200$ polynomial pairs at each of the
three fixed scales $1$, $2$ and $\tfrac32$, that is $600$ scale--pair cases,
with seed $1729$; the non-integer fixed scale is not sampled by the first
harness in the same way. Each case checks addition, multiplication, and strict
reversal of the support order under $a\mapsto-ca$. The output is recorded in
\texttt{results/verification.json}; all $257$ exact coefficients of both $B$
and $H$, together with the coefficients of $B^2$ and $H^2$, are in
\texttt{results/coefficients.csv}.

These computations do not verify infinite summability, the surreal number
axioms, or the universal quantifiers of the theorems. The all-degrees proofs
are the coefficient recursion in \cref{lem:witness} and the derivative
argument in \cref{lem:H}. The checks instead detect sign, indexing,
coefficient, and finite-convolution errors in the concrete formulas and
implementation.

\subsection{A compact logical audit}
The contradiction uses an element that really belongs to the proposed source:
its support is the allowed increasing sequence $0,1,2,\ldots$, and its
coefficients are nonzero real numbers. The identity is proved in that source
using its own multiplication. Unitality follows from the prescribed image of
$x^0$. The target inequality is a field-order statement, not an ordinal
calculation. No image is assigned to an infinite sum by distributing $\Phi$
over its terms. These checks isolate all assumptions needed for the answer.

\subsection{What the result does not settle}
The conclusions above are specific to \emph{multiplicative} evaluations on the
stated series domains. They do not prohibit useful regularizations of selected
divergent series, they say nothing about nonmultiplicative summation
conventions, and they do not constrain game-valued operations formulated with
different axioms. They likewise do not constitute a classification of
arbitrary surreal field embeddings. As recorded in \cref{sec:claims},
Lipparini's broader Problem 7.5 remains open as far as this note is concerned:
a partial operation, an operation on a smaller domain, or an operation not
required to be multiplicative evades everything proved here.

\section{Conclusion}
The requested extension cannot exist. The full obstruction is already
\[
 \boxed{\quad B(x)^2=1-x\quad\text{in }\Q[[x]],\qquad 1-\omega<0
        \quad\text{in }\No.\quad}
\]
The increasing-support convention carries an algebraic commitment: its
positive exponent monomials behave as infinitesimals in any ordered-field
image. This commitment can be read from square roots of units even when no
order-preservation axiom is imposed.

The obstruction is already countable, rational and quadratic. It persists
without continuity, without order preservation, without any hypothesis about
infinite sums, and --- by \cref{thm:semiring} --- even in a
nonnegative-coefficient semiring formulation in which subtraction is
unavailable in the source and the excluded values are widened from $u>1$ to
$u\ge1$.

There is nevertheless an exact positive picture. Ordinary power series can
be evaluated at every surreal infinitesimal, and no other surreal values of
the variable are possible. General pure monomial evaluations exist precisely
for additive, strictly order-reversing exponent maps. The canonical example
is $x^a\mapsto\omega^{-a}$. Thus the failure at $\omega$ is not a shortage of
surreal values or a class-size difficulty: it is a reversal of the magnitude
orientation required by the source algebra.

\appendix
\section{Well-ordering and summability from finite words}\label{app:summability}
This appendix proves the support facts used in the extensions. It is not
needed for the main contradiction, whose coefficients have finite ordinary
Cauchy convolutions. We work in the usual set-theoretic setting with choice.
Every support below is a set, even if the ambient ordered group is a class.

\subsection{Well-ordered sets and finite products}\label{app:products}
A set in a linear order is \emph{well ordered} when every nonempty subset has
a least element, and \emph{reverse well ordered} when every nonempty subset
has a greatest element. The two notions are not interchangeable:
$\{0,1,2,\ldots\}$ is well ordered and $\{0,-1,-2,\ldots\}$ is reverse well
ordered, and neither set has the other property. The whole orientation
question of this note is the difference between them.

\begin{lemma}[Nondecreasing subsequences]\label{lem:subseq}
Every infinite sequence in a well-ordered set has an infinite nondecreasing
subsequence. Consequently, every infinite sequence in a finite product of
well-ordered sets has an infinite coordinatewise nondecreasing subsequence.
\end{lemma}
\begin{proof}
For the first claim, choose an occurrence of the least value taken by the
whole sequence, then an occurrence of the least value taken on the remaining
infinite tail, and continue. The successive tail minima cannot decrease, since
each is the minimum of a subset of the previous tail.

For the second, apply the first claim to the first coordinate to obtain an
infinite subsequence nondecreasing there, then to the second coordinate of
that subsequence, and so on. A finite product requires only finitely many
extractions, and each extraction preserves the monotonicity already achieved.
\end{proof}

\subsection{Finite words over a well-ordered alphabet}
Let $A$ be a well-ordered set. For finite words
$u=(u_1,\ldots,u_m)$ and $v=(v_1,\ldots,v_n)$ over $A$, write $u\preceq v$
when there are indices
\[
 1\le i_1<\cdots<i_m\le n
 \quad\text{such that}\quad u_j\le v_{i_j}\ (1\le j\le m).
\]
The empty word embeds into every word. A sequence is called \emph{bad} if no
earlier member embeds into a later one.

\begin{lemma}[Higman's lemma, well-ordered alphabet case]\label{lem:higman}
Every infinite sequence of finite words over $A$ contains an increasing pair
for $\preceq$.
\end{lemma}
\begin{proof}
Suppose a bad infinite sequence exists. Construct a bad sequence
$w_0,w_1,\ldots$ with the following minimality property: having chosen
$w_0,\ldots,w_{n-1}$, choose $w_n$ to have the smallest possible length among
words permitting that prefix to be extended to a bad infinite sequence.
The possible lengths form a nonempty subset of $\N$, so a least one exists
at every stage.

No $w_n$ is empty, since an empty word embeds into every later word. Write
$w_n=a_n v_n$, with first letter $a_n\in A$ and a shorter tail $v_n$.
By the first part of \cref{lem:subseq}, choose
\[
 i_0<i_1<i_2<\cdots,
 \qquad a_{i_0}\le a_{i_1}\le a_{i_2}\le\cdots.
\]
Consider the modified sequence
\[
 w_0,\ldots,w_{i_0-1},v_{i_0},v_{i_1},v_{i_2},\ldots.
\]
It is bad. An embedding between two words in the unchanged prefix would
already contradict badness of the original sequence. If a prefix word
embeds into a tail $v_{i_j}$, it embeds into $w_{i_j}$ as well, again a
contradiction. Finally, if $v_{i_j}\preceq v_{i_k}$ for $j<k$, adjoining the
first letters with $a_{i_j}\le a_{i_k}$ gives
$w_{i_j}\preceq w_{i_k}$, also impossible.

The modified sequence therefore extends the prefix through index $i_0-1$
to a bad sequence using the strictly shorter word $v_{i_0}$ at the next
position. This contradicts the chosen minimality of $w_{i_0}$.
\end{proof}

This is the standard minimal-bad-sequence mechanism for the special case of
Higman's theorem \citep{Higman1952} used below; no general tree theorem is needed.

\subsection{The positive-support lemma}
\begin{lemma}[Positive-support summability]\label{lem:neumann}
Let $G$ be a linearly ordered abelian group, and let $A\subseteq G_{>0}$ be
well-ordered. Let $M(A)$ be the set of sums of finite words over $A$, including
the empty sum $0$. Then:
\begin{enumerate}[label=(\alph*)]
\item $M(A)$ is well-ordered.
\item For every $g\in G$, only finitely many finite words over $A$ have sum $g$.
\end{enumerate}
In (b) it makes no difference whether words are counted as nondecreasing words
or as arbitrary ordered finite tuples: a fixed finite word has only finitely
many permutations, so the two counts are finite together.
\end{lemma}
\begin{proof}
If $u\preceq v$, the sum of the letters of $u$ is at most the sum of the
letters of $v$. If $u\ne v$, this inequality is strict: either there is an
extra positive letter in $v$, or one of the matched letters is strictly
larger. When the word lengths agree, the only strictly increasing bijection
of their index sets is the identity, so the latter alternative exhausts
the distinct-word case.

If $M(A)$ contained an infinite strictly decreasing sequence of sums, choose
one representing word for each sum. An increasing pair of words provided by
\cref{lem:higman} would give a nondecreasing pair of sums, a contradiction.
A linearly ordered set with no infinite strictly decreasing sequence is
well-ordered in the present set-theoretic setting: from a nonempty subset
with no minimum one could choose a descending sequence. This proves (a).

If infinitely many distinct words had the same sum $g$, enumerate a countable
sequence of distinct such words. An increasing pair from
\cref{lem:higman} would have strictly increasing sums by the preceding
observation, again a contradiction. This proves (b).
\end{proof}

Notice that no Archimedean property is assumed. The positive elements of $A$
may themselves be infinitesimal in a larger ordered field, and $A$ may have
any set-sized ordinal order type.

The hypothesis $A\subseteq G_{>0}$ is not a convenience. If $0$ were allowed
into $A$, then part (b) would fail immediately: appending any number of zeros
to a word representing $g$ produces another word representing $g$, so every
element of $M(A)$ would have infinitely many representations, and the
regrouping arguments that use (b) would collapse. Part (a) would survive, but
it is (b) that the substitution results need.

\subsection{Finite convolution and formal substitution}
\begin{corollary}\label{cor:convolution}
If $A,B\subseteq G$ are well-ordered, then $A+B$ is well-ordered, and for each
$c\in G$ only finitely many pairs $(a,b)\in A\times B$ satisfy $a+b=c$.
\end{corollary}
\begin{proof}
There is nothing to prove if either set is empty. Otherwise let $a_0,b_0$
be their minima. The set
\[
 C=\bigl((A-a_0)\cup(B-b_0)\bigr)\setminus\{0\}
\]
is well-ordered and consists of positive elements. A finite union of
well-ordered subsets of a linear order is well-ordered, since the minimum of
any nonempty subset is the smaller of its existing minima in the two parts.

Now $A+B\subseteq a_0+b_0+M(C)$, proving the first claim by
\cref{lem:neumann}. For the second, translate a pair to its nonnegative
components and remove any zero entries. It becomes a word of length at most
two over $C$ with sum $c-a_0-b_0$. There are finitely many such words, and each
word has only finitely many ways to restore the two labeled positions and
possible zeros. Thus there are finitely many pairs.
\end{proof}

The two-series case deserves a second, independent proof, because it is the
only support fact used by the definition of the Hahn product itself. The
following argument needs neither Higman's lemma nor the positive-support
monoid, only \cref{lem:subseq}. Keeping it makes the \emph{definition} of
$\Hu{G}$ independent of any well-quasi-order machinery; the
well-quasi-order machinery is needed only for infinite substitutions.

\begin{lemma}[Two supports, elementarily]\label{lem:two-supports}
Let $A,B$ be well-ordered subsets of a linearly ordered abelian group $G$.
Then $A+B$ is well ordered, and every $c\in G$ has only finitely many
representations $c=a+b$ with $(a,b)\in A\times B$.
\end{lemma}
\begin{proof}
Suppose $c_0>c_1>c_2>\cdots$ were an infinite strictly decreasing sequence in
$A+B$, and write $c_n=a_n+b_n$ with $a_n\in A$ and $b_n\in B$. By the second
part of \cref{lem:subseq} applied to $A\times B$, pass to a subsequence along
which both $a_n$ and $b_n$ are nondecreasing. Then $c_n$ is nondecreasing
along that subsequence, contradicting strict decrease. In a linear order, the
absence of an infinite strictly decreasing sequence gives well-ordering in the
present set-theoretic setting: from a nonempty subset with no least element
one recursively chooses a descending sequence.

For the second claim, suppose some $c$ had infinitely many distinct
representing pairs $(a_n,b_n)$. Again extract a subsequence that is
coordinatewise nondecreasing. Two distinct pairs $(a_i,b_i)$ and $(a_j,b_j)$
in it satisfy $a_i\le a_j$ and $b_i\le b_j$ with at least one inequality
strict, so $a_i+b_i<a_j+b_j$; but both sums equal $c$.
\end{proof}

\begin{corollary}[Substitution into a positive-support series]\label{cor:substitution}
If $\epsilon\in\Hu{G}$ has support contained in $G_{>0}$, then for any real
sequence $(d_n)_{n\ge0}$ the Hahn sum
\[
 \sum_{n\ge0}d_n\epsilon^n
\]
is well-defined. Sending a formal power series $\sum d_nz^n$ to this sum
is a unital homomorphism $\R[[z]]\to\Hu{G}$.
\end{corollary}
\begin{proof}
For $\epsilon=0$ this is the constant-term map. Otherwise apply
\cref{lem:neumann} to $A=\supp(\epsilon)$. Every exponent occurring in a
power lies in $M(A)$, and each exponent can occur in only finitely many
powers because it has finitely many word representations, across all word
lengths. Thus the family is summable.

To multiply two resulting sums, expand products as finite words in $A$.
At a specified exponent there are finitely many words and finitely many
splits of each word into two parts. Consequently the product can be regrouped
by the two lengths, giving the ordinary formal power-series product.
Preservation of addition and the unit is immediate.
\end{proof}

This corollary justifies the binomial expressions in
\cref{lem:positivesquare} without an analytic convergence argument. It also
explains why setting the formal variable equal to an infinitesimal is natural:
positive source exponent supports correspond, after reflection, to negative
surreal normal-form exponents.

\subsection{The field property}
\begin{proposition}\label{prop:field}
For any linearly ordered abelian group $G$, $\Hu{G}$ is a field.
\end{proposition}
\begin{proof}
By \cref{cor:convolution}, addition and multiplication are well-defined.
Their ring identities follow from finite regrouping of coefficients; for
three or more factors, iterate that corollary to see that each coefficient
still has only finitely many contributing tuples.

Let $s\ne0$ have leading term $r x^a$, where $r\ne0$. Write
$s=r x^a(1+\epsilon)$ with $\supp(\epsilon)\subseteq G_{>0}$. By
\cref{cor:substitution}, the formal geometric identity gives
\[
 (1+\epsilon)\sum_{n\ge0}(-\epsilon)^n=1.
\]
Thus
\[
 s^{-1}=r^{-1}x^{-a}\sum_{n\ge0}(-\epsilon)^n
\]
is a well-defined inverse. All nonzero elements are invertible.
\end{proof}

\section{Reproducing the artifact checks}\label{app:reproduce}
The archive is self-contained apart from a standard LaTeX installation and
Python 3.9 or later. It has no network-dependent verification steps and no
third-party Python package requirements.

From the extracted archive directory, run
\begin{lstlisting}
python code/verify.py --degree 256 --trials 1000 --output results
\end{lstlisting}
A successful run prints a JSON report with \texttt{"status": "PASS"} and
rewrites the exact coefficient table and verification report. The explicit
runtime checks are written with a \texttt{require} helper that raises a
dedicated \texttt{VerificationFailure} exception, so they remain active even
when Python is invoked with \texttt{-O}, and a failure exits with a nonzero
status after printing a \texttt{FAIL} line rather than a traceback. The
optional flags are \texttt{--degree}, \texttt{--trials},
\texttt{--fixed-pairs}, \texttt{--seed} and \texttt{--output}; the first three
reject nonpositive values.

The file \texttt{results/coefficients.csv} has one header row and one row per
degree $0\le n\le256$, with the seven columns \texttt{n},
\texttt{numerator}, \texttt{denominator} and \texttt{exact\_fraction} for
$b_n$, then \texttt{h\_n}, \texttt{coefficient\_B\_squared} and
\texttt{coefficient\_H\_squared}. The last two columns are the computed
coefficients of $B^2$ and $H^2$ and should read $1,-1,0,0,\ldots$ and
$1,1,1,\ldots$ respectively.

\begin{center}
\begin{tabular}{p{0.30\textwidth}p{0.60\textwidth}}
\toprule
Artifact & Purpose\\
\midrule
\texttt{article.tex}, \texttt{article.pdf}
 & Standalone source and compiled article.\\
\texttt{short-proof.tex}, \texttt{short-proof.pdf}
 & One-page extraction of the decisive argument.\\
\texttt{code/verify.py}
 & Exact coefficient and finite-exponent checks.\\
\texttt{results/verification.json}
 & Machine-readable verification results and parameters.\\
\texttt{results/coefficients.csv}
 & Exact coefficients of both explicit witness series.\\
\texttt{references.bib}
 & Reusable bibliographic metadata; not a build dependency.\\
\texttt{proof-audit.md}, \texttt{sources.md}
 & Logical audit, and version-specific source record.\\
\texttt{README.md}, \texttt{build.sh}
 & Reproduction instructions and a local build command.\\
\bottomrule
\end{tabular}
\end{center}

Build the article using
\begin{lstlisting}
latexmk -pdf -interaction=nonstopmode -halt-on-error article.tex
\end{lstlisting}
or the supplied \texttt{build.sh}. The document uses standard TeX Live
packages and the bundled installation's Palatino-style \texttt{newpx} fonts;
no font files are redistributed. The typeset bibliography is included directly in the LaTeX source, so no
BibTeX run is required. Reusable bibliographic metadata are also supplied in
\texttt{references.bib}.

The supplied \texttt{proof-audit.md} records the logical assumptions and the
limits of the computational evidence. \texttt{sources.md} records inspected
primary sources and the status-search boundary.

\clearpage
\phantomsection
\addcontentsline{toc}{section}{References}
\begin{thebibliography}{9}
\bibitem[Berarducci and Mantova(2018)]{BerarducciMantova2018}
Alessandro Berarducci and Vincenzo Mantova.
\newblock Surreal numbers, derivations and transseries.
\newblock \emph{Journal of the European Mathematical Society}, 20:339--390, 2018.
\newblock \href{https://doi.org/10.4171/JEMS/769}{doi:10.4171/JEMS/769}.
The inspected preprint is \href{https://arxiv.org/abs/1503.00315v3}{arXiv:1503.00315v3};
see Sections 2.3--2.6.

\bibitem[van den Dries and Ehrlich(2001)]{DriesEhrlich2001}
Lou van den Dries and Philip Ehrlich.
\newblock Fields of surreal numbers and exponentiation.
\newblock \emph{Fundamenta Mathematicae}, 167(2):173--188, 2001.
\newblock \href{https://doi.org/10.4064/fm167-2-3}{doi:10.4064/fm167-2-3}.
The exponent-reversing embedding is displayed on p.~176.

\bibitem[Kaplan et al.(2025)]{KKS}
Elliot Kaplan, Lothar Sebastian Krapp, and Michele Serra.
\newblock Decomposing the automorphism group of the surreal numbers.
\newblock \href{https://arxiv.org/abs/2509.22374v1}{arXiv:2509.22374v1},
26 September 2025.
\newblock Cited for standard normal-form and Hahn-operation background, and for
the distinction between arbitrary and strongly additive maps; see Sections
1--4. This preprint was carried over from the merged manuscript and was not
independently re-inspected for the present note.

\bibitem[Bournez and Guilmant(2022)]{BG}
Olivier Bournez and Quentin Guilmant.
\newblock Surreal fields stable under exponential and logarithmic functions.
\newblock \href{https://arxiv.org/abs/2201.08199}{arXiv:2201.08199}, 2022.
\newblock Cited for standard omega-map background only. This preprint was
carried over from the merged manuscript and was not independently
re-inspected for the present note.

\bibitem[Higman(1952)]{Higman1952}
Graham Higman.
\newblock Ordering by divisibility in abstract algebras.
\newblock \emph{Proceedings of the London Mathematical Society},
series 3, 2(1):326--336, 1952.
\newblock \href{https://doi.org/10.1112/plms/s3-2.1.326}{doi:10.1112/plms/s3-2.1.326}.

\bibitem[Lipparini(2026)]{Lipparini2026}
Paolo Lipparini.
\newblock Monotone infinitary operations on ordinals (extended version).
\newblock \href{https://arxiv.org/abs/2505.00424v2}{arXiv:2505.00424v2},
30 April 2026.
\newblock Problem 7.7, p.~38. The earlier formulation is Problem 6.8 in
version 1, dated 1 May 2025.
\end{thebibliography}
\end{document}
